\chapter{Phương pháp}
\label{chap:phuong-phap}

\section{Câu hỏi nghiên cứu và giả thuyết}
\label{sec:rq}

\textbf{RQ1 (tái hiện):} kết quả công bố của SemLoc có tái hiện được không, và còn đứng vững khi chấm công bằng, khi bỏ
điểm thưởng không được mô tả trong bài báo và khi kiểm tra lại nhãn (dòng lỗi chuẩn) bằng thực thi không?
\textbf{RQ2 (LLM cần gì):} một prompt duy nhất có kiểm thử thất bại có định vị đúng ít nhất bằng SemLoc không, với chi phí
bao nhiêu; tín hiệu nào tạo ra khác biệt (kiểm thử thất bại hay độ phủ); và các cách bổ sung (CovGate, kiểm chứng bằng
bản sửa, biểu quyết, phân tích tĩnh) có thêm được gì không?
\textbf{RQ3 (độ bền):} kết luận có còn đúng trên lỗi thật (BugsInPy), trên mã đã đổi tên định danh,
với năm họ mô hình và trên cả ba bộ nhãn (gốc, đã chuẩn hóa, đã rà soát) không?
Câu hỏi chung, thứ hạng của SemLoc đến từ bằng chứng thực thi hay dòng LLM tự đoán, được kiểm qua giả thuyết H1.

\textbf{H1:} ở hạng 1, độ chính xác của SemLoc đến từ dòng LLM tự đoán, bằng chứng thực thi chỉ góp ở top-3 (bác bỏ nếu
riêng phổ vi phạm đạt Acc@1 không thấp hơn SemLoc). \textbf{H2:} prompt có kiểm thử thất bại không kém SemLoc (bác bỏ nếu
khoảng tin cậy của chênh lệch Acc@1 hoặc Acc@3 nằm hẳn dưới 0). \textbf{H3:} thêm độ phủ từng dòng vào prompt cải thiện
kết quả (bác bỏ nếu khoảng tin cậy không nằm hẳn trên 0).

\section{Dữ liệu}
\label{sec:du-lieu}

\textbf{SemFault-250}~\cite{yang2026semloc} gồm 250 chương trình Python một hàm, mỗi chương trình có một lỗi cài vào và
khoảng mười kiểm thử; nhóm chia cố định 50 chương trình làm tập dev (chỉ để chọn tham số) và 200 làm tập test. Trên nhãn
đã chuẩn hóa, tập test còn \RqTwoRepairedAllSemlocN{} chương trình vì các nhãn không kiểm chứng được bị loại.

\textbf{Bộ đã chuẩn hóa.} Có \GtBlankN{} dòng lỗi chuẩn là dòng trống và \GtNotExecutedN{} chương trình không có kiểm
thử thất bại nào chạy tới dòng lỗi chuẩn. Nhóm xây một \emph{bộ phán định sửa tối thiểu} (Hình~\ref{fig:sua-nhan}): một
dòng là vị trí lỗi nếu một thay đổi nhỏ tại dòng đó làm \emph{mọi} kiểm thử đạt. LLM chỉ đề xuất bản sửa, kiểm thử quyết
định (GPT-6 luna đề xuất; các lượt chạy lại dùng Claude Sonnet 5.5 và Claude Opus 5.5); dòng lân cận chỉ được tìm khi
dòng gốc không chạy tới được. Bản sửa ``nhỏ'' thay một dòng và đổi tối đa bốn token hoặc một phần ba dòng. Chỉ bản thay thế
mới xác nhận được một câu lệnh; điền dòng trống hay chèn thêm một dòng (cách sửa lỗi thiếu mã) chỉ được chấp nhận khi không
câu lệnh nào gần đó sửa được, vì dòng điền hay chèn thêm thường chỉ ghi đè câu lệnh lỗi ở chỗ khác. Kết quả: \GtConfirmedN{} nhãn được xác nhận,
\GtCorrectedN{} được dời vị trí, \GtUnverifiedN{} chưa kiểm chứng được (giữ vị trí gốc) và \GtUnverifiableN{} bị loại;
bộ nhãn đã chuẩn hóa còn \GtRepairedN{} chương trình.

\begin{figure}[H]\centering
% Bộ phán định sửa tối thiểu: từ nhãn gốc đến bộ nhãn đã sửa. Số liệu qua macro.
{\setstretch{1}%
\begin{tikzpicture}[
  q/.style={fbox=semloc, text width=46mm, minimum height=12mm},
  leaf/.style={fbox=#1, text width=33mm, minimum height=15mm}, leaf/.default=gray,
  lab/.style={font=\small, fill=white, inner sep=1.5pt}]
\node[q] (q1) at (0, 0) {Có kiểm thử thất bại chạy tới nhãn gốc?};
\node[q] (q2) at (-4.1, -2.0) {Bản sửa nhỏ \emph{tại nhãn}\\làm mọi kiểm thử đạt?};
\node[q] (q3) at (4.1, -2.0) {Bản sửa nhỏ ở dòng lân cận\\($\le 4$ dòng) làm mọi kiểm thử đạt?};
\node[leaf=failloc] (l1) at (-6.15, -4.7) {\textbf{Xác nhận}\\\GtConfirmedN{} nhãn};
\node[leaf] (l2) at (-2.05, -4.7) {\textbf{Chưa kiểm chứng}\\\GtUnverifiedN{} nhãn\\(giữ nhãn gốc)};
\node[leaf=failloc] (l3) at (2.05, -4.7) {\textbf{Dời vị trí}\\\GtCorrectedN{} nhãn};
\node[leaf] (l4) at (6.15, -4.7) {\textbf{Không kiểm chứng}\\\textbf{được}: \GtUnverifiableN{}\\(loại)};
\draw[farr] (q1.south) -- ++(0,-3mm) -| node[lab, pos=0.25, above] {có} (q2.north);
\draw[farr] (q1.south) -- ++(0,-3mm) -| node[lab, pos=0.25, above] {không} (q3.north);
\draw[farr] (q2.south) -- ++(0,-3mm) -| node[lab, pos=0.25, above] {có} (l1.north);
\draw[farr] (q2.south) -- ++(0,-3mm) -| node[lab, pos=0.25, above] {không} (l2.north);
\draw[farr] (q3.south) -- ++(0,-3mm) -| node[lab, pos=0.25, above] {có} (l3.north);
\draw[farr] (q3.south) -- ++(0,-3mm) -| node[lab, pos=0.25, above] {không} (l4.north);
\node[fbox=failloc, text width=120mm] (out) at (-1.0, -7.0)
  {\textbf{SemFault-250 đã sửa: \GtRepairedN{} chương trình}\\\small bỏ nhãn không kiểm chứng được và
   \GtBlankFixDroppedN{} vị trí sửa nằm trên dòng trống};
\foreach \l in {l1, l2, l3} \draw[farr] (\l.south) -- (\l.south |- out.north);
\end{tikzpicture}}

\caption{Bộ phán định sửa tối thiểu: LLM chỉ đề xuất bản sửa, kiểm thử quyết định dòng nào là dòng lỗi. Nhãn được xác nhận hoặc dời vị trí khi một bản sửa nhỏ làm mọi kiểm thử đạt; nhãn chưa kiểm chứng giữ nguyên vị trí gốc}\label{fig:sua-nhan}
\end{figure}

\textbf{Bộ đã rà soát và duyệt mù.} Một LLM (Claude Fable 5.1) đọc chương trình, kiểm thử và nhãn gốc của cả 250
chương trình và nhận định nhãn đúng, sai (kèm dòng đề xuất) hay lỗi hai dòng. \ReviewReviewersN{} thành viên nhóm duyệt
\ReviewHumanN{} nhãn (các nhãn đã dời và mọi nhãn LLM cho là sai), đồng ý với LLM ở \ReviewAgreeN{}/\ReviewHumanN{} ca;
trưởng nhóm sửa thêm \ReviewOverrideN{} ca có bằng chứng chạy kiểm thử; \LabaudAllExtended{} nhãn được thêm dòng thứ hai
chỉ theo LLM. Bộ \emph{đã rà soát} có đủ \LabaudAllN{} chương trình, \LabaudAllTwoLine{} nhãn hai dòng. Vòng \emph{duyệt
mù} trên \BlindBlindAllPrograms{} chương trình ngẫu nhiên kiểm tra lại bộ này: hai thành viên chưa từng xem chương trình
duyệt độc lập, dòng ứng viên được xáo và không kèm nhận định của LLM; \BlindAnswersExcluded{}/\BlindAnswersTotal{} câu
dẫn lại tệp duyệt cũ bị loại theo quy tắc đặt trước. Bộ đã chuẩn hóa chỉ do kiểm thử quyết định nên dùng cho kết quả
chính; bộ đã rà soát có nhãn do LLM đề xuất nên có thể có lợi cho phương pháp dùng LLM (Mục~\ref{sec:moi-de-doa}) và
chỉ là góc nhìn độ nhạy (Bảng~\ref{tab:ba-bo-nhan}, Hình~\ref{fig:ra-soat}).

\begin{table}[H]\centering
\caption{Ba bộ dòng lỗi chuẩn dùng trong tiểu luận}\label{tab:ba-bo-nhan}
\small
\begin{tabularx}{\linewidth}{@{}l>{\raggedright\arraybackslash}X>{\raggedright\arraybackslash}X@{}}
\toprule Bộ nhãn & Ai quyết định & Vai trò \\ \midrule
Gốc & Tác giả SemLoc & So với bài báo gốc \\
Đã chuẩn hóa & Kiểm thử (LLM chỉ đề xuất bản sửa) & Kết quả chính; LLM không phán đoán vị trí \\
Đã rà soát & LLM rà trước, thành viên nhóm duyệt, trưởng nhóm sửa ca có bằng chứng chạy kiểm thử & Góc nhìn độ nhạy; đủ mọi chương trình \\
\bottomrule\end{tabularx}
\end{table}

\begin{figure}[H]\centering
% Rà soát nhãn hai bước và duyệt mù. Số liệu qua macro.
{\setstretch{1}%
\begin{tikzpicture}[node distance=6mm and 5mm,
  s/.style={fbox=semloc, text width=33mm, minimum height=19mm},
  f/.style={fbox=failloc, text width=33mm, minimum height=19mm},
  g/.style={fbox=gray, text width=33mm, minimum height=19mm}]
\node[g] (a) {\textbf{Nhãn gốc}\\\LabaudAllN{} chương trình};
\node[s, right=of a] (b) {\textbf{LLM rà trước}\\đọc mã, test, nhãn;\\nhận định từng nhãn};
\node[s, right=of b] (c) {\textbf{Người duyệt}\\\ReviewReviewersN{} thành viên duyệt\\\ReviewHumanN{} nhãn nghi vấn};
\node[f, right=of c] (d) {\textbf{Nhãn đã rà soát}\\\LabaudAllMoved{} nhãn dời,\\\LabaudAllTwoLine{} nhãn hai dòng};
\node[s, text width=88mm, anchor=north east] (e) at ($(d.south west)+(-8mm,-12mm)$) {\textbf{Duyệt mù}: \BlindBlindAllPrograms{} chương trình chọn ngẫu
  nhiên; mỗi chương trình do hai thành viên chưa từng xem duyệt độc lập; các dòng ứng viên được đánh chữ theo thứ tự ngẫu
  nhiên, không ghi nguồn, không có gợi ý của LLM};
\node[f, text width=40mm, anchor=north] (r) at ($(d.south)+(0,-12mm)$) {\textbf{Kết quả}: \BlindBlindAllMatchReviewed/\BlindBlindAllNReviewed{} câu trả lời
  trùng nhãn đã rà soát, \BlindBlindAllMatchOriginal/\BlindBlindAllNOriginal{} trùng nhãn gốc};
\draw[farr] (a) -- (b); \draw[farr] (b) -- (c); \draw[farr] (c) -- (d);
\draw[farr] (d.south) -- (d.south |- r.north); \draw[farr] (e) -- (r);
\end{tikzpicture}}

\caption{Rà soát nhãn hai bước (LLM rà trước, thành viên nhóm duyệt) và vòng duyệt mù kiểm tra lại bộ nhãn đã rà soát mà không có gợi ý của LLM}\label{fig:ra-soat}
\end{figure}

\textbf{BugsInPy và biến thể đổi tên.} Từ dự án youtube-dl của BugsInPy~\cite{widyasari2020}, nhóm giữ 31 lỗi tái hiện
được (kiểm thử thất bại trên bản lỗi, đạt trên bản sửa, đo được độ phủ); dòng lỗi chuẩn là các dòng bị sửa trong hàm mà
kiểm thử thất bại chạy qua. Để kiểm tra ghi nhớ~\cite{ramos2024memorize,xu2024contam}, mọi định danh tự định nghĩa được
đổi thành tên vô nghĩa và chú thích bị xóa, giữ nguyên số dòng và kết quả kiểm thử: còn 199 chương trình của tập test.

\section{Các phương pháp so sánh}
\label{sec:phuong-phap-so-sanh}

Bảng~\ref{tab:phuong-phap} tóm tắt năm phương pháp của so sánh chính, Hình~\ref{fig:thiet-ke} là thiết kế thực nghiệm.
LLM-Direct là chế độ dự đoán trực tiếp của SemLoc; SemLoc chạy bằng mã nguồn gốc của tác giả, chỉ thay lớp gọi LLM.

\begin{table}[H]\centering
\caption{Các phương pháp trong so sánh chính: đầu vào, số lượt gọi LLM trung bình mỗi chương trình (tập test, Gemini 3.8 Flash) và việc có chạy kiểm thử hay không}\label{tab:phuong-phap}
{\setstretch{1.15}\small
\begin{tabularx}{\linewidth}{@{}lXcc@{}}
\toprule Phương pháp & Đầu vào & Lượt gọi LLM & Chạy kiểm thử \\ \midrule
SBFL (Ochiai) & Ma trận độ phủ kiểm thử $\times$ dòng & -- & có \\
LLM-Direct & Chỉ mã nguồn & \CostLlmDirectCalls & không \\
SemLoc & Mã nguồn; ràng buộc do LLM viết được chèn vào và chạy cùng kiểm thử & \CostSemlocCalls & có \\
SemLoc−Bonus & Như SemLoc, bỏ điểm thưởng cho dòng LLM tự đoán & \CostSemlocCalls & có \\
\textbf{FailLoc} & Mã nguồn, tên và dòng lỗi đầu tiên của từng kiểm thử thất bại & \CostLlmTestsCalls & có \\
\bottomrule\end{tabularx}}
\end{table}

\begin{figure}[H]\centering
% Thiết kế thực nghiệm lần 1: dữ liệu × phương pháp × bộ nhãn × lần chạy.
{\setstretch{1}%
\begin{tikzpicture}[node distance=4mm and 6mm,
  col/.style={fbox=#1, text width=34mm, minimum height=44mm, anchor=north},
  head/.style={font=\bfseries, anchor=south}]
\node[col=gray] (d) at (0,0) {SemFault-250\\[2mm] tập dev: 50 CT\\(chọn tham số)\\[1mm] tập test: 200 CT\\(báo cáo)\\[2mm]\small chia phân tầng,\\hạt giống cố định};
\node[col=semloc, right=of d.north east, anchor=north west] (m) {SBFL (Ochiai)\\LLM-Direct\\SemLoc\\SemLoc−Bonus\\\textcolor{failloc}{\textbf{FailLoc}}\\[2mm]\small cùng mô hình\\Gemini 3.8 Flash};
\node[col=failloc, right=of m.north east, anchor=north west] (g) {Nhãn gốc\\[1mm] Nhãn đã sửa\\bằng thực thi\\[2mm]\small \GtRepairedN{} CT\\sau khi sửa};
\node[col=gray, right=of g.north east, anchor=north west] (r) {3 lần chạy\\mỗi cấu hình\\[2mm] Acc@1, Acc@3,\\Acc@5\\[2mm]\small bootstrap theo\\chương trình và\\theo cụm mẫu};
\node[head] at (d.north) {Dữ liệu};
\node[head] at (m.north) {Phương pháp};
\node[head] at (g.north) {Bộ nhãn};
\node[head] at (r.north) {Lặp và chỉ số};
\foreach \a/\b in {d/m, m/g, g/r} \node at ($(\a.east)!0.5!(\b.west)$) {\Large$\times$};
\end{tikzpicture}}

\caption{Thiết kế thực nghiệm của so sánh chính: mọi phương pháp chạy trên cùng các chương trình của tập test, cùng mô hình, trên ba bộ nhãn (gốc, đã chuẩn hóa, đã rà soát), mỗi cấu hình có LLM chạy ba lần}\label{fig:thiet-ke}
\end{figure}

\textbf{FailLoc} (Failure-guided Localization), phương pháp đề xuất của nhóm, là một lượt gọi LLM
($T = \ConstTemperature$) với mã nguồn đánh số dòng, số kiểm thử đạt/không đạt, câu ``chương trình có đúng một dòng
lỗi'', tên và dòng thông báo lỗi đầu tiên của mỗi kiểm thử thất bại; LLM trả về năm dòng nghi ngờ nhất
(Hình~\ref{fig:failloc-vd}; prompt trong kho mã, tệp \texttt{src/flx/direct\_cov.py}, cố định trước khi xem tập test).
Ý tưởng này lấy từ thực tế công nghiệp~\cite{medeiros2026crash} và từ kết quả nhật ký lỗi giúp LLM định vị~\cite{wu2023llmfl}. Biến thể \textbf{FailLoc+Cov} chú thích thêm cho mỗi dòng số kiểm thử thất bại và thành công
chạy qua nó: FailLoc so với LLM-Direct đo tác dụng của kiểm thử thất bại, FailLoc+Cov so với FailLoc đo tác dụng của độ phủ.

\begin{figure}[H]\centering
{\small\setstretch{1.0}
\fbox{\parbox{0.95\linewidth}{\ttfamily\footnotesize The program below has exactly one faulty line. 1 tests fail and 4 pass.\\
~\\
...\\
9: def votes\_apportion(total, weights, *, floor\_to\_int=True, minimum=0):\\
10: ~~~~"""\\
11: ~~~~Compute weighted allocation for votes capacity.\\
12: ~~~~"""\\
13: ~~~~if len(weights) == 0:\\
14: ~~~~~~~~raise ValueError("weights required")\\
15: ~~~~if total < 0:\\
16: ~~~~~~~~raise ValueError("negative total")\\
17: ~~~~if minimum < 0:\\
18: ~~~~~~~~raise ValueError("negative minimum")\\
19: ~~~~if sum(weights) == 0:\\
20: ~~~~~~~~raise ValueError("zero total weight")\\
21: \\
22: ~~~~weight\_sum = sum(weights)\\
23: ~~~~shares = []\\
24: ~~~~for w in weights:\\
25: ~~~~~~~~portion = (w / weight\_sum) * total\\
26: ~~~~~~~~shares.append(portion if portion > minimum else minimum)\\
27: \\
28: ~~~~if floor\_to\_int:\\
29: ~~~~~~~~return [int(v) for v in shares]\\
30: ~~~~return shares\\
~\\
Failing tests:\\
- test\_truncation\_sum\_equals\_total\_fail: \\
~\\
Return the five most suspicious line numbers, most suspicious first.}}

\medskip
\fbox{\parbox{0.95\linewidth}{Câu trả lời của LLM (năm dòng nghi nhất): \texttt{[29, 25, 26, 28, 30]}; dòng lỗi theo nhãn: 29.}}}

\caption{FailLoc trên \texttt{votes\_apportion} (lượt chạy 1, Gemini 3.8 Flash): prompt gửi cho LLM (các dòng ngoài hàm được lược bớt) và năm dòng LLM trả về; dòng lỗi theo nhãn đứng hạng 1}\label{fig:failloc-vd}
\end{figure}

\textbf{Các cách bổ sung.} \emph{CovGate} bổ sung bằng chứng độ phủ cho điểm dòng của SemLoc qua bốn thành phần bật/tắt
được (cổng độ phủ, cộng $\lambda\cdot$Ochiai với $\lambda \in \{\ConstLambdaGrid\}$, xếp lại bằng kiểm chứng phản thực
tế, chỉ thưởng dòng đoán khi có bằng chứng thực thi); cấu hình chọn trên tập dev chỉ còn là phá hòa bằng Ochiai
($\lambda = \ConstLambdaOriginal$ trên bộ gốc, $\lambda = \ConstLambdaRepaired$ trên bộ đã chuẩn hóa).
\emph{FailLoc+Verify} (Hình~\ref{fig:verify}) dùng kiểm thử làm trọng tài: với $k$ dòng đầu của FailLoc, LLM viết một bản
sửa nhỏ tại dòng đó, và dòng có bản sửa làm mọi kiểm thử đạt được đưa lên đầu; lưới \PavGridSize{} cấu hình được đánh
giá offline trên tập dev với quy tắc chọn cố định trước. \emph{FailLoc+Vote} gộp ba câu trả lời bằng tổng nghịch đảo hạng.

\begin{figure}[H]\centering
% FailLoc+Verify: đề xuất rồi kiểm chứng bằng bản sửa và test. Số liệu qua macro.
{\setstretch{1}%
\begin{tikzpicture}[node distance=6mm and 5mm,
  f/.style={fbox=failloc, text width=29mm, minimum height=16mm},
  s/.style={fbox=semloc, text width=29mm, minimum height=16mm},
  g/.style={fbox=gray, text width=29mm, minimum height=16mm}]
\node[f] (a) {\textbf{FailLoc đề xuất}\\năm dòng nghi nhất};
\node[s, right=of a] (b) {\textbf{LLM viết bản sửa nhỏ}\\cho từng dòng\\trong $k$ dòng đầu};
\node[s, right=of b] (c) {\textbf{Chạy bộ test}\\trên chương trình\\đã vá};
\node[f, right=of c] (d) {\textbf{Xếp lại}\\dòng có bản sửa\\làm mọi test đạt\\lên trước};
\draw[farr] (a) -- (b); \draw[farr] (b) -- (c); \draw[farr] (c) -- (d);
\node[g, below=9mm of b, text width=66mm, xshift=17mm] (p) {\textbf{Vì sao không cải thiện}: test của bộ dữ liệu quá yếu
  để làm trọng tài. Trên các dòng ứng viên \emph{không} có lỗi, bản sửa nhỏ vẫn làm mọi test đạt ở
  \PavmechReviewedOtherLineFixable\% lần thử; tại dòng lỗi thật chỉ \PavmechReviewedFaultLineFixable\% (tập dev, nhãn đã rà soát)};
\draw[farr=gray] (c.south) -- (c.south |- p.north);
\end{tikzpicture}}

\caption{FailLoc+Verify (đề xuất rồi kiểm chứng): kiểm thử được dùng làm trọng tài cho vị trí lỗi; trên bộ dữ liệu này kiểm thử quá yếu nên bản sửa nhỏ ở dòng không lỗi cũng thường làm mọi kiểm thử đạt}\label{fig:verify}
\end{figure}

\textbf{Phân tích tĩnh.} Ruff và pyrefly (phiên bản cố định) chạy trên cả 250 chương trình; một chương trình ``trúng''
khi có cảnh báo đúng dòng lỗi chuẩn, và số trúng được so với số trúng kỳ vọng nếu cảnh báo rải ngẫu nhiên.

\section{Kiểm tra cách đánh giá và phân rã đóng góp}
\label{sec:audit-phan-ra}

Từ nhật ký công bố của tác giả, nhóm tính lại mọi cấu hình theo \emph{cách chấm của tác giả} (phá hòa theo thứ tự dòng,
chương trình chạy lỗi bị loại khỏi mẫu số) và \emph{cách chấm công bằng} (công thức~\eqref{eq:acck}), rồi bỏ lần lượt
các thành phần không được mô tả. Để phân rã ``grounded hay guessed'', giá trị Shapley~\eqref{eq:shapley} chia Acc@$k$
của SemLoc cho hai tín hiệu \{dòng đoán, phổ\} (bản mở rộng thêm Ochiai và kiểm chứng phản thực tế).

\section{Thiết lập thực nghiệm}
\label{sec:thiet-lap}

Mô hình chính là Gemini 3.8 Flash; RQ3 thêm GPT-6 luna, Claude Sonnet 5.5, gpt-oss-120b và Gemini 3.5 Flash-Lite trên
một tập con cố định 50 chương trình của tập test (\ParetoSemlocGeminiThreeEightFlashN{} trên nhãn đã chuẩn hóa). SemLoc chạy ở $T = \ConstTemperature$ như bài báo; mọi lượt gọi LLM được ghi nhật ký
và lưu đệm nên phát lại được. Mỗi cấu hình có LLM chạy 3 lần, chỉ dùng lượt qua kiểm tra sức khỏe (không còn lỗi gọi
LLM hay lỗi ở bước kiểm chứng phản thực tế). Kết quả là trung bình $\pm$ độ lệch chuẩn qua 3 lần; so sánh cặp dùng
$\Delta$ kèm khoảng tin cậy bootstrap theo cụm, McNemar và Wilcoxon (mục~\ref{sec:thong-ke}). Chi phí đo bằng số lượt
gọi, số token và thời gian mỗi chương trình.
